\documentclass[11pt,letterpaper]{article}
\usepackage[margin=0.8in]{geometry}
\usepackage[T1]{fontenc}
\usepackage{lmodern,amsmath,amssymb,microtype}
\usepackage[colorlinks=true,urlcolor=blue,linkcolor=blue]{hyperref}
\setlength{\parindent}{0pt}
\setlength{\parskip}{5pt}
\newcommand{\nystrom}{Nystr\"om}
\hypersetup{pdftitle={Residual-guided Nystrom preconditioning: initial experiments}}
\begin{document}
\begin{center}
{\large\bfseries Residual-guided \nystrom{} preconditioning: initial experiments}\\[4pt]
CS 6220 research record \qquad November 23, 2026
\end{center}

\textbf{Question and implemented method.}
Can the current residual guide kernel-column replacements that reduce the total
time to reach a prescribed solution accuracy, while keeping the number of
retained columns fixed?
For the Gaussian kernel matrix $K$, the implementation solves
$A\alpha=y$, where $A=K+\lambda I$ and $\lambda>0$.
Randomly pivoted Cholesky (RPCholesky) selects an initial set $S$ of $k$ columns. The
preconditioner and exchange score are
\[
M_S=K_{:,S}K_{S,S}^{\dagger}K_{S,:}+\lambda I,
\qquad q(S;r)=r^{\mathsf T}M_S^{-1}r,
\qquad r=y-A\alpha.
\]
Here $K_{:,S}$ contains the selected columns, $K_{S,S}$ contains their entries
at the selected rows, and $\dagger$ denotes the pseudoinverse.
Every 20 iterations of preconditioned conjugate gradients (PCG), four candidate
columns are sampled sequentially with probabilities proportional to the squared
residual entries outside $S$. Each candidate is compared with all possible
removals, including discarding the candidate itself. A replacement is accepted
only if a stable evaluation of $q$ decreases by more than $10^{-12}$ relative.
The solver restarts from its current iterate after an accepted checkpoint and
continues its recurrence after a rejected checkpoint. Exactly $k$ distinct
column indices are retained throughout.

\textbf{Exchange score and subsequent progress.}
The \nystrom{} approximation satisfies
$K_{:,S}K_{S,S}^{\dagger}K_{S,:}\preceq K$, where $B\preceq C$ means
$v^{\mathsf T}Bv\leq v^{\mathsf T}Cv$ for every vector $v$. Consequently,
\[
e^{\mathsf T}Ae=r^{\mathsf T}A^{-1}r\leq q(S;r),
\qquad e=\alpha-A^{-1}y.
\]
An accepted exchange tightens this upper bound without changing the current solution
error. Whether it improves the next iterations is an empirical question.
The instructor's suggested paper,
\href{https://arxiv.org/abs/2509.21963v2}{\emph{IterativeCUR: Large Rank-Adaptive
Approximation From a Small Recycled Sketch}}, uses the matrix residual of a CUR
approximation to select additional rows and columns. Here the residual belongs
to a linear solve, and each accepted new column replaces a retained column.
The paper motivates reusing matrix information during selection; the following
experiment directly tests the effect of the proposed fixed-budget exchanges.

\textbf{Comparing preconditioners from the same iterate.}
For each of five seeds, synthetic inputs contain four Gaussian clusters and are
paired with a smooth target and a Gaussian target independent of the inputs. The two targets
share the same kernel and data split. There are 512 training observations,
$k=16$, and $\lambda\in\{10^{-3},10^{-5}\}$.
At the first three accepted checkpoints of each solve, two diagnostic branches
start from the same iterate: one restarts with the old preconditioner, and one
restarts with the exchanged preconditioner. Each preconditioner is then held
fixed for up to 20 iterations. A direct Cholesky solution supplies the reference
for the squared energy error $e^{\mathsf T}Ae$.

\begin{center}
\small
\begin{tabular}{lrrrrr}
\hline
$\lambda$ & Checkpoints & $q_{\rm new}/q_{\rm old}$ & $E^{\rm new}_5/E^{\rm old}_5$ & $E^{\rm new}_{10}/E^{\rm old}_{10}$ & $E^{\rm new}_{20}/E^{\rm old}_{20}$ \\
\hline
$10^{-3}$ & 30 & 0.955 & 1.135 & 1.282 & 1.024 \\
$10^{-5}$ & 30 & 0.955 & 1.004 & 1.005 & 1.001 \\
\hline
\end{tabular}
\end{center}
Here $E_h=e_h^{\mathsf T}Ae_h$ after $h$ branch iterations. Entries are medians
of checkpoint ratios; an energy ratio above one favors the old preconditioner.


At all 60 inspected checkpoints, the updated set reduced $q$. Nevertheless, the exchanged
preconditioner produced a larger squared energy error after both 5 and 10
iterations at every inspected checkpoint. After 20 iterations it produced a
smaller error at 24 checkpoints and a larger error at 36. The checkpoint
observations are grouped within 20 solves; they are not independent replicates.
At $\lambda=10^{-3}$, all ten adaptive solves reached the relative residual
tolerance $10^{-8}$. At $\lambda=10^{-5}$, none reached it within 2,000 iterations.

\textbf{Total solve time.}
The timing comparison uses the two synthetic targets, UCI
\href{https://archive.ics.uci.edu/dataset/165/concrete+compressive+strength}{Concrete Compressive Strength},
and UCI \href{https://archive.ics.uci.edu/dataset/291/airfoil+self+noise}{Airfoil Self-Noise}.
Each case uses an 80/20 train/test split,
training-only standardization, and a Gaussian bandwidth equal to the median
training pairwise distance. The settings are $k=16$, $\lambda=10^{-3}$, five
seeds, and three timing repeats per seed. The initial selected columns are
shared across methods. ``Fixed'' holds the original preconditioner throughout;
``Restarted fixed'' uses those same columns and restarts at the accepted-update
iterations of the adaptive solve; ``Exchange'' is the residual-guided method.
Unpreconditioned conjugate gradients (CG) is also included.

Each solve uses one CPU BLAS thread on an Intel Core i7-13700H under WSL2 and
the same blockwise kernel operator.
Timing includes initialization, updates, and directly checked residuals.
Dense reference construction, preprocessing, and post-solve prediction checks
are outside the common timing boundary. The stopping rule requires a true
relative residual at most $10^{-8}$, with a 2,000-iteration limit.
Method order is randomized for each repeat; a separate untimed adaptive run
establishes the restart schedule, which is checked against the timed run.

\begin{center}
\small\setlength{\tabcolsep}{4pt}
\begin{tabular}{lrrrrr}
\hline
Workload & CG (s) & Fixed (s) & Restarted fixed (s) & Exchange (s) & Time ratio \\
\hline
Synthetic (smooth) & 1.031 & 0.397 & 0.837 & 0.916 & 2.24 \\
Synthetic (random) & 1.111 & 0.420 & 0.927 & 0.967 & 2.36 \\
Concrete & 4.409 & 1.976 & 3.724 & 5.615 & 2.89 \\
Airfoil & 7.264 & 2.511 & 4.438 & 5.871 & 2.34 \\
\hline
\end{tabular}
\end{center}
Times are medians across the five seeds of each seed's median over three repeats.
The final column is the median of the five paired exchange/fixed time ratios.


All 240 completed timed solves reached the stopping tolerance. The exchange method was slower than fixed PCG in 20 of the 20 paired dataset--target--seed cases. Restarted fixed PCG required more iterations than uninterrupted fixed PCG in 20 cases; exchanging columns required more iterations than restarted fixed PCG in 19 cases. Exchange-update work accounted for a median 5.0\% of adaptive solve time across cases. The additional iterations therefore deserve attention alongside update cost.

The largest measured true relative residual was $9.98\times 10^{-9}$. The measured maxima of relative coefficient error and relative prediction difference from the direct reference were $6.90\times 10^{-8}$ and $2.16\times 10^{-8}$, respectively. These maxima cover all methods, seeds, and repeats in this comparison.

The first execution group reached its 900-second wall-time allowance after 237 completed measurements. The remaining three measurements were completed separately with the same source, settings, and prepared data. The interrupted attempt and incomplete record remain preserved.


An additional one-seed execution check ran all eight implemented methods on
the smooth synthetic case and Airfoil. In addition to the four methods above,
it included target-based exchanges, uniform candidate sampling, periodic
RPCholesky reconstruction, and limited-memory BFGS using block secant pairs.
All 16 solves converged. Their repeated multi-seed comparison, the larger column
budgets, and the remaining regularization settings belong to the next stages
of the original research plan.

\newpage
\textbf{Next experiment.}
The branch diagnostic already compares old and new columns under the same
restart. A third branch will continue the old PCG recurrence without restarting,
preserving its recursive residual, search direction, and recurrence scalars. Comparing this branch
with restarted old PCG will measure the effect of restarting at the same
iterate; comparing restarted old and new PCG will measure the effect of the
column exchange.

The next diagnostic will also record the actual decrease available in one
restarted PCG step. For $z=M^{-1}r$, the step length is
$(r^{\mathsf T}z)/(z^{\mathsf T}Az)$ and the squared energy decrease is
\[
\Delta(M;r)=\frac{(r^{\mathsf T}z)^2}{z^{\mathsf T}Az}.
\]
In the independently checked first checkpoint of the smooth seed-0 case at
$\lambda=10^{-3}$, exchanging columns reduces $q$ from $11132.41$ to $10796.33$,
but reduces $\Delta$ from $76.45$ to $73.25$. The direction and its curvature
therefore explain a distinction that the upper-bound score misses.
Checking $\Delta$ requires an additional matrix--vector product for each
preconditioner evaluated. Any later use of this quantity for accepting exchanges
will include that cost in total solve time. The research question remains the
fixed-column-budget improvement of the solve; the present evidence calls for
revising the exchange criterion and examining the restart policy.

\textbf{Reproduction.}
The exact configurations are in \texttt{configs/}; complete histories, exchange
decisions, source snapshots, and version records are in \texttt{results/}.
\texttt{results/initial-analysis.json} contains the paired summaries used here.
\texttt{docs/branch-verification.txt} records the independent numerical check.
The implementation passes 19 numerical and solver tests. The submitted proposal
is preserved in \texttt{proposal/}.
\end{document}
